\subsection{Restrições de Atuadores (Torque/Força)}

Além dos limites cinemáticos e geométricos das juntas, o manipulador robótico está sujeito às restrições dinâmicas existentes nos atuadores responsáveis por gerar o movimento.  
Essas restrições são expressas em termos de torque máximo nas juntas revolutas e força máxima na junta prismática.

A seguir, as restrições são expressas formalmente para cada tipo de junta:

\begin{equation}
\label{eq:limite_torque}
|\tau_j| \leq \tau_{j,\text{max}} \qquad \text{(juntas revolutas)},
\end{equation}

\begin{equation}
\label{eq:limite_forca}
|f_5| \leq f_{5,\text{max}} \qquad \text{(junta prismática)}.
\end{equation}

Esses valores máximos foram definidos de forma arbitrária - utilizando como base as simulações realizadas - e para esta dissertação, os valores de $\tau_{j,\text{max}}$ e $f_{5,\text{max}}$ serão determinados posteriormente com base na análise dinâmica completa e na sintonia do controlador PD implementado no ambiente \emph{Simulink}.  

